\documentclass[11pt]{article}
\usepackage[margin=1in]{geometry}
\usepackage{amsmath,amssymb}
\usepackage{booktabs}
\usepackage{array}
\usepackage[table]{xcolor}
\usepackage{hyperref}
\hypersetup{colorlinks=true,linkcolor=blue,urlcolor=blue}
\usepackage{parskip}

\newcommand{\xor}{\oplus}
\newcommand{\code}[1]{\texttt{#1}}
\newcommand{\maj}{\mathrm{maj}}

\title{\bfseries Gadget Leakage Investigation}
\author{}
\date{2026-08-25}

\begin{document}
\maketitle
\vspace{-3em}

\noindent\textbf{Scope.} From the gadgetization tool through a cut-adversary
analysis of a production deliverable (\code{run\_c}). Covers: (I) the trace leak in
the single-carrier gadgetizer; (II) the new nonlinear \code{193}/\code{291} gadgets
and why they resist the linear attack; (III) a gauntlet re-run on the production
single-carrier path; and (IV) a segment-relation (``cut adversary'') study of
\code{run\_c}'s phase-A and final outputs, including dense-local-window equation
recovery and the effect of phase~B. All results are empirical under the stated
threat models, not proofs; single run (\code{run\_c}), degree-1 relations except
where noted.

\section{Context: the gadgetizer and the trace leak}

The production path is the \textbf{single-carrier swap-refresh} gadgetizer
(\code{ProdConfig::production\_single()}). It is \textbf{additive product-sharing}:
each logical value $v$ lives on one carrier as $\text{carrier} = v \xor M(\text{band})
\xor \kappa$, and every monomial of the mask $M$ is XORed onto the carrier by a gate.
Under the gauntlet's \textbf{\code{xtrace}} attack --- a full-trace affine adversary
observing the initial wires plus \emph{every gate flip} --- the payload is recovered
exactly,
\[
v = \text{carrier}(t) \;\xor\; \Big(\bigoplus_s f_s\Big) \;\xor\; (\text{linear band
corrections}) \;\xor\; (\kappa\text{-discharge flips}),
\]
every term of which is a trace column. This is \textbf{structural and
$n$-independent}: the mask monomials are \emph{materialized as flips}, so scale does
not help. The design study named three candidate fixes (withheld/never-materialized
masks; a nonlinear multi-carrier decode; certified block fusion). The
\code{193}/\code{291} gadgets are the second lever.

\section{The nonlinear gadgets (\code{193}, \code{291})}

\textbf{Encoding.} Each logical bit is stored in a five-wire block via a
\emph{nonlinear} (degree-2) decode
\[
E(x_0,x_1,x_2,x_3,x_4) = x_0 \xor x_1 \xor \maj(x_2,x_3,x_4), \qquad
\maj = x_2x_3 \xor x_3x_4 \xor x_2x_4,
\]
with \textbf{two shares} per mutable logical wire ($v = E(s_1)\xor E(s_2)$). Cost:
\textbf{10 physical wires per logical bit} --- a single r57 gate uses \textbf{44
wires} (\code{193}) or \textbf{71 wires} (\code{291}), $\approx 5\times$ the two-wire
additive share.

\textbf{How it computes without materializing the secret.} The gate re-shares the
target onto fresh wires. Correctness reduces to producing an output-majority triple
whose majority equals the old majority XOR the gate value XOR a pad. This is done
with a \textbf{secret-controlled majority switch} (\code{SB\_U0} identity /
\code{SB\_U1} majority-flipping) and a \textbf{masked build $\to$ run $\to$ unbuild}
sequence. The key discipline: an \textbf{operand product is always ANDed with the
switch or a blinding mask, never emitted bare}; masks cancel only in the logical
sum, and chaff terms cancel pairwise, so no single flip equals a secret product.
\code{193} emits fan-in-$\le 4$; \code{291} is a \emph{mask-first fan-in-2 lowering}
of the same logical stream ($+27$ shared decomposition ancillas, a persistent CSE
cache), with identical security by construction.

\textbf{Linear-relations check.} Enumerating the exact GF(2) relations among
$\{\,\text{const},\ \text{init wires},\ \text{all gate flips}\,\}$ (held-out verified):
\begin{itemize}
\item All secrets --- $a$, $b$, $a\!\cdot\!b$, $c_{\text{in}}$, $c_{\text{out}}$ ---
are \textbf{not} in the affine span, per gate \emph{and} across a 2-gate chain.
\item The only relations are harmless: ancilla-restoration telescopes, duplicate-flip
cancellations (chaff / build--unbuild), and copies of the \emph{linear pad} wires.
Every relation touches only linear components; the nonlinear $\maj$ core (hence every
secret) never enters a linear relation.
\end{itemize}
So \code{193}/\code{291} resist \code{xtrace} at the tested scale --- the opposite of
the additive path. (Verified on single gates and a 2-gate chain; production scale
untested.)

\section{Gauntlet on the production single-carrier gadgetizer}

The gauntlet's \code{bandproduct92} arm \emph{is} \code{production\_single()},
exercised at $n=8$ on a source chain the harness fabricates (so it knows the
ground-truth targets); it does not ingest production artifacts. Re-run (fresh
binaries, mixing off):

\begin{center}
\begin{tabular}{r l l}
\toprule
$k$ & \code{xtrace} (single-carrier) & \code{xtrace} (\code{none} control) \\
\midrule
1 & 2/2 & 2 \\
2 & 4/4 & 4 \\
16 & \textbf{64/64} & 64 \\
\bottomrule
\end{tabular}
\end{center}

\noindent\code{xtrace} is \textbf{fully saturated at $k=16$ --- identical to the
unprotected control}. Snapshot attacks \emph{do} drop (\code{a1} $64\!\to\!23$,
\code{xrows} $48\!\to\!16$): swap-refresh protects cuts, not the full trace. This
reproduces the prior result and confirms the additive-path leak is structural.

\section{\texttt{run\_c} cut-adversary analysis (\texttt{segment\_deduce})}

\textbf{Instrument.} \code{segment\_deduce} drives the \textbf{source} (the sandwich
containing C) and a \textbf{predictor} (a mixed circuit) on the same input $x$ and
asks, for each source \emph{segment} (a wire's value between two writes), whether it
is a low-degree GF(2) function of the predictor's \textbf{wire values at $K$ cuts
(register snapshots)}. This is the \textbf{cut adversary} --- snapshots only, \emph{no
gate flips} --- strictly weaker than \code{xtrace}. A ``recovered'' C-segment equals a
fixed XOR of observed snapshot values.

Source $=$ \code{gss.mpmct1.sandwich.mpmct1} (256 wires, 7{,}920 gates); predictor $=$
\code{phaseA.mpmct1} (512 wires, 1{,}410{,}438 gates). Segment budget: 8{,}176 total
($7{,}920$ writes $+ 256$ initial values); region \textbf{C} $= 2{,}680$ (the
strictly-pre-neck slice; $\text{C}+\text{D}+\text{band} = 6{,}000 = 2\times \code{mcd
3000}$).

\subsection{Global dense-cut sweep (phase-A)}

\begin{center}
\begin{tabular}{r l r}
\toprule
pred-cuts & region-C recovered & interior-only leaks (all regions) \\
\midrule
8   & $7/2680 = 0.3\%$   & 0 \\
32  & $7/2680 = 0.3\%$   & --- \\
64  & $28/2680 = 1.0\%$  & 76 \\
128 & $137/2680 = 5.1\%$ & 754 \\
\bottomrule
\end{tabular}
\end{center}

Region-C recoverability \textbf{grows with cut resolution}. The 8-cut ``recoveries''
are all \emph{trivial output-boundary} equalities $C(w) = \code{phaseA}(w{+}128)$ at
the output cut (the sandwich's own $y{=}x$ neck copy); \textbf{zero genuine interior
leaks} at 8 cuts. Genuine interior leaks appear at $64+$ cuts: multi-term,
\textbf{always cross-wire} (never same-wire), spanning 1--36 cuts. Cut spacing: 64
cuts $\approx 22$k gates/cut, 128 $\approx 11$k.

\subsection{Dense local window (the efficient attack)}

One \textbf{20k-gate window} (phase-A gates $400{,}000$--$420{,}000$, $1.4\%$ of the
circuit) with \textbf{128 cuts placed inside it} (156 gates/cut --- $\sim\!70\times$
finer than the global scan):
\begin{itemize}
\item \textbf{162 C-values recovered} --- \emph{more than the entire global 128-cut
scan} (137).
\item \textbf{Temporally localized:} all recovered C-values are born in source gates
\textbf{1{,}420--2{,}293}; D/S2/N $=0$ in this window.
\end{itemize}
So C's intermediates are \textbf{temporally localized}, dense local cutting is far
more efficient than global cutting, and the global scan \emph{understated} the leak
(it was resolution-limited).

\subsection{Anatomy of the 162 equations}
\begin{itemize}
\item \textbf{Terms per equation:} median 7, mean 9.7, range 3--64.
\item \textbf{Spread:} median gate-span $\approx 10{,}080$ (half the window); median
wire-index range 363 of 512.
\item \textbf{Coverage:} 101 of C's 128 logical wires; 358 of 512 predictor wires used.
\item \textbf{Patterns:} (i) \emph{same-wire telescoping} --- one wire read at many
times, XORed; (ii) \emph{hub wires} recurring across equations (w90 in 30, w410 in 25,
w32 in 24); (iii) a shared \emph{entry-snapshot} anchor; (iv) genuine \emph{cross-wire}
structure --- a carrier's history XORed with the band wires that masked it.
\end{itemize}

\subsection{What is special about the hub wires}

The hubs are \textbf{the most-written wires inside the window} --- w419 (\#4 of 512),
w90 (\#9), w410 (\#10), w32 (\#13) --- \emph{not} globally special. Mechanism: many
writes $\Rightarrow$ many distinct values across cuts $\Rightarrow$ a rich
reconstruction basis; read-heavy but write-light wires are \emph{not} hubs.
Structurally they split into \textbf{data/carrier wires} (w90, w32, w31)
\textbf{$+$ band/mask wires} (w410, w511, w302, w419) --- exactly the
$\text{carrier}\xor\text{mask}$ pair. This is also why dense \emph{local} cutting
works: in the window these carrier/band wires change value dozens of times.

\subsection{What phase B adds (final vs phase-A)}

Equivalent dense window on \code{final.mpmct1} (post split/cross/fcompress,
$4{,}859{,}911$ gates), matched resolution (156 gates/cut), comparable C-content
(source $\sim\!1500$--$2200$):

\begin{center}
\begin{tabular}{l r r}
\toprule
 & phase-A window & final window \\
\midrule
C-values recovered      & 162 & 46 \\
source-content covered  & 873 gates & 655 gates \\
\textbf{C-leak density} & \textbf{0.186 / src-gate} & \textbf{0.070 / src-gate} \\
C-wires hit             & 101 / 128 & 45 / 128 \\
top hub frequency       & w90 in 30 eqns & w53 in 6 eqns \\
\bottomrule
\end{tabular}
\end{center}

Phase B adds \textbf{genuine hiding, $\approx 2.7\times$} (density $0.186\to 0.070$,
beyond the length-dilution confound), and \textbf{de-concentrates} the leak (the hub
structure collapses; the dominant phase-A hubs w90/w410 vanish, only w355 survives in
both). It \textbf{does not close} the leak: one final window still recovers
\textbf{45 of 128 C-wires}.

\section{Takeaways}
\begin{enumerate}
\item \textbf{Additive single-carrier is broken against the trace adversary}
(\code{xtrace} 64/64), structurally and at any scale.
\item \textbf{Against the weaker cut adversary}, the same construction holds at coarse
snapshots but \textbf{leaks under dense/local cutting}, growing with resolution. The
leak is \textbf{temporally localized} and concentrates on the \textbf{actively-written
carrier/band wires} of each window.
\item \textbf{Phase B roughly triples the local hiding and breaks up the hub
structure}, but does not eliminate the cut-leak.
\item \textbf{The nonlinear \code{193}/\code{291} encoding keeps every secret out of
the linear span} (per-gate and 2-gate chain) --- it targets the \code{xtrace} leak
directly, at $\approx 5\times$ the wire cost. Production-scale evaluation is the open
item.
\end{enumerate}

\noindent\textbf{Caveats.} One run (\code{run\_c}); degree-1 relations except the
nonlinear check; single windows (not a full sweep); nonlinear analysis at small
scale; the cut adversary is \emph{not} the trace adversary.

\noindent\textbf{Tooling.} \code{segment\_deduce} was extended with
\code{--cut-lo}/\code{--cut-hi} (place cuts in a gate window) and \code{--eq-dump}
(dump full region-C equations); the original source was backed up.

\end{document}
